\documentclass[12pt,a4paper]{article}
\usepackage[utf8]{inputenc}
\usepackage[T1]{fontenc}
\usepackage{amsmath}
\usepackage{amsfonts}
\usepackage{amssymb}
\usepackage{graphicx}
\usepackage{hyperref}
\usepackage{listings}
\usepackage{xcolor}
\usepackage{geometry}
\geometry{margin=1in}

\title{Infinity Mode Protocol Extension\\RFC 25.02202\\Dynamic Asset Loading Protocol}
\author{Tootsville Development Team}
\date{\today}

\begin{document}

\maketitle

\begin{abstract}
This RFC proposes an extension to the Infinity Mode protocol to formalize the dynamic loading of 3D assets and game resources. The extension establishes standardized mechanisms for requesting, caching, and managing game assets from remote servers, ensuring efficient resource utilization while maintaining protocol compliance. This addresses the need to move from hard-coded asset references to server-driven asset management in the modern 3D continuous world environment.
\end{abstract}

\section{Introduction}

The current Tootsville implementation uses hard-coded asset references and geometries for 3D models and game resources. This RFC proposes a protocol extension that enables:

\begin{itemize}
\item Dynamic asset discovery and loading
\item Server-driven asset management
\item Efficient caching and resource optimization
\item Progressive loading for improved performance
\item Protocol-compliant asset distribution
\end{itemize}

\section{Asset Loading Architecture}

\subsection{Current Limitations}

The existing implementation suffers from:

\begin{itemize}
\item Hard-coded model URIs and geometries
\item Static asset bundles with no dynamic loading
\item Inefficient resource utilization
\item Lack of server control over asset distribution
\item No progressive loading capabilities
\end{itemize}

\subsection{Proposed Architecture}

The new architecture provides:

\begin{itemize}
\item Server-driven asset discovery
\item Dynamic loading based on scene requirements
\item Intelligent caching with LRU eviction
\item Progressive loading with quality levels
\item Fallback mechanisms for missing assets
\end{itemize}

\section{Protocol Extension Specification}

\subsection{Asset Request Message Format}

Asset requests SHALL use the following JSON format:

\begin{lstlisting}[language=JSON, caption=Asset Request Message Format]
{
  "type": "asset-request",
  "request-id": "unique-request-identifier",
  "timestamp": "ISO-8601-timestamp",
  "assets": [
    {
      "asset-id": "unique-asset-identifier",
      "type": "model|texture|sound|animation",
      "format": "glb|gltf|png|jpg|mp3|wav",
      "quality": "low|medium|high",
      "priority": "critical|normal|optional"
    }
  ],
  "context": {
    "scene-id": "current-scene-identifier",
    "player-position": {
      "x": 0.0,
      "y": 0.0,
      "z": 0.0
    },
    "viewport": {
      "width": 1920,
      "height": 1080,
      "fov": 75
    }
  }
}
\end{lstlisting}

\subsection{Asset Response Message Format}

Asset responses SHALL use the following format:

\begin{lstlisting}[language=JSON, caption=Asset Response Message Format]
{
  "type": "asset-response",
  "request-id": "original-request-identifier",
  "timestamp": "ISO-8601-timestamp",
  "status": "success|partial|error",
  "assets": [
    {
      "asset-id": "unique-asset-identifier",
      "status": "available|loading|error",
      "uri": "https://assets.tootsville.org/models/asset.glb",
      "metadata": {
        "size": 1024000,
        "checksum": "sha256-hash",
        "version": "1.0.0",
        "dependencies": ["texture-id-1", "animation-id-1"]
      },
      "quality": "high",
      "priority": "critical"
    }
  ],
  "cache-control": {
    "max-age": 3600,
    "etag": "asset-etag",
    "last-modified": "ISO-8601-timestamp"
  }
}
\end{lstlisting}

\subsection{Asset Types and Formats}

\subsubsection{3D Models}

\begin{itemize}
\item \textbf{Primary Format}: GLTF/GLB 2.0
\item \textbf{Fallback Formats}: OBJ, FBX (converted to GLTF)
\item \textbf{Compression}: Draco geometry compression for GLTF
\item \textbf{Lodding}: Automatic level-of-detail generation
\end{itemize}

\subsubsection{Textures}

\begin{itemize}
\item \textbf{Formats}: WebP, PNG, JPEG
\item \textbf{Compression}: Basis Universal for GPU textures
\item \textbf{Mipmaps}: Automatic mipmap generation
\item \textbf{Resolution}: Adaptive based on distance and quality settings
\end{itemize}

\subsubsection{Sounds}

\begin{itemize}
\item \textbf{Formats}: Opus, MP3, WAV
\item \textbf{Spatial}: 3D positional audio support
\item \textbf{Streaming}: Progressive loading for long audio files
\end{itemize}

\subsubsection{Animations}

\begin{itemize}
\item \textbf{Format}: GLTF animations
\item \textbf{Blending}: Support for animation blending
\item \textbf{Procedural}: Runtime animation generation
\end{itemize}

\section{Asset Discovery and Loading}

\subsection{Scene-Based Discovery}

Assets SHALL be discovered based on scene descriptions:

\begin{lstlisting}[language=JSON, caption=Scene Description with Assets]
{
  "scene": {
    "id": "butterfly-frame-area",
    "bounds": {
      "min": [-100, -100, -100],
      "max": [100, 100, 100]
    },
    "objects": [
      {
        "id": "butterfly-frame-1",
        "type": "scenery",
        "position": [10, 0, 5],
        "model-uri": "/Assets/Models/5/ButterflyFrame/ButterflyFrame.glb",
        "quality": "high"
      }
    ]
  }
}
\end{lstlisting}

\subsection{Progressive Loading}

Assets SHALL be loaded progressively based on:

\begin{enumerate}
\item \textbf{Critical Assets}: Required for scene functionality
\item \textbf{Normal Assets}: Standard quality assets
\item \textbf{Optional Assets}: Enhanced quality or decorative assets
\end{enumerate}

\subsection{Caching Strategy}

\subsubsection{Client-Side Caching}

Clients SHALL implement a multi-tier caching system:

\begin{itemize}
\item \textbf{Memory Cache}: Recently used assets
\item \textbf{Disk Cache}: Persistent storage with LRU eviction
\item \textbf{Service Worker Cache}: Offline capability
\end{itemize}

\subsubsection{Cache Management}

Cache SHALL be managed using:

\begin{itemize}
\item LRU (Least Recently Used) eviction policy
\item Size-based limits (per asset type)
\item Version-based invalidation
\item Network-aware preloading
\end{itemize}

\section{Implementation Guidelines}

\subsection{Client Implementation}

\subsubsection{Vue Component for Dynamic Asset Loading}

\begin{lstlisting}[language=TypeScript, caption=DynamicAssetLoader Component]
import { GLTFLoader } from 'three/examples/jsm/loaders/GLTFLoader.js';
import { DRACOLoader } from 'three/examples/jsm/loaders/DRACOLoader.js';

interface AssetRequest {
  assetId: string;
  type: 'model' | 'texture' | 'sound';
  uri: string;
  priority: 'critical' | 'normal' | 'optional';
  quality: 'low' | 'medium' | 'high';
}

class DynamicAssetLoader {
  private loader: GLTFLoader;
  private dracoLoader: DRACOLoader;
  private cache: Map<string, any> = new Map();
  private loadingQueue: AssetRequest[] = [];

  constructor() {
    this.loader = new GLTFLoader();
    this.dracoLoader = new DRACOLoader();
    this.dracoLoader.setDecoderPath('/draco/');
    this.loader.setDRACOLoader(this.dracoLoader);
  }

  public async loadAsset(request: AssetRequest): Promise<any> {
    // Check cache first
    if (this.cache.has(request.assetId)) {
      return this.cache.get(request.assetId);
    }

    // Add to loading queue
    this.loadingQueue.push(request);

    // Process queue based on priority
    return this.processQueue();
  }

  private async processQueue(): Promise<any> {
    // Sort by priority
    this.loadingQueue.sort((a, b) => {
      const priorityOrder = { critical: 0, normal: 1, optional: 2 };
      return priorityOrder[a.priority] - priorityOrder[b.priority];
    });

    // Load highest priority asset
    const request = this.loadingQueue.shift();
    if (!request) return null;

    try {
      const asset = await this.performLoad(request);
      this.cache.set(request.assetId, asset);
      return asset;
    } catch (error) {
      console.error(`Failed to load asset ${request.assetId}:`, error);
      return this.getFallbackAsset(request);
    }
  }

  private async performLoad(request: AssetRequest): Promise<any> {
    switch (request.type) {
      case 'model':
        return new Promise((resolve, reject) => {
          this.loader.load(
            request.uri,
            (gltf) => resolve(gltf),
            (progress) => {
              // Handle loading progress
            },
            reject
          );
        });
      // Handle other asset types...
    }
  }

  private getFallbackAsset(request: AssetRequest): any {
    // Return appropriate fallback based on asset type
    return null;
  }
}
\end{lstlisting}

\subsubsection{Asset Manager Integration}

The asset loader SHALL integrate with the main game engine:

\begin{lstlisting}[language=TypeScript, caption=Asset Manager Integration]
class AssetManager {
  private loader: DynamicAssetLoader;
  private sceneManager: SceneManager;
  private networkManager: NetworkManager;

  public async loadSceneAssets(sceneId: string): Promise<void> {
    // Request asset list from server
    const assetList = await this.networkManager.requestSceneAssets(sceneId);

    // Load critical assets first
    const criticalAssets = assetList.filter(a => a.priority === 'critical');
    await Promise.all(criticalAssets.map(a => this.loader.loadAsset(a)));

    // Load remaining assets in background
    const remainingAssets = assetList.filter(a => a.priority !== 'critical');
    remainingAssets.forEach(asset => {
      this.loader.loadAsset(asset).catch(console.error);
    });
  }

  public getAsset(assetId: string): any {
    return this.loader.getAsset(assetId);
  }
}
\end{lstlisting}

\subsection{Server Implementation}

\subsubsection{Asset Registry}

The server SHALL maintain an asset registry:

\begin{lstlisting}[language=Lisp, caption=Asset Registry Implementation]
(defclass asset-registry ()
  ((assets :initform (make-hash-table :test 'equal))
   (metadata-cache :initform (make-hash-table :test 'equal))))

(defmethod register-asset ((registry asset-registry) asset-info)
  "Register an asset in the registry"
  (let ((asset-id (getf asset-info :id)))
    (setf (gethash asset-id (assets registry)) asset-info)
    (update-metadata-cache registry asset-id)))

(defmethod get-asset-info ((registry asset-registry) asset-id)
  "Retrieve asset information"
  (gethash asset-id (assets registry)))

(defmethod get-assets-for-scene ((registry asset-registry) scene-id)
  "Get all assets required for a scene"
  (let ((scene-assets (gethash scene-id (metadata-cache registry))))
    (mapcar (lambda (asset-id) (get-asset-info registry asset-id))
            scene-assets)))
\end{lstlisting}

\subsubsection{Asset Serving}

Assets SHALL be served with appropriate HTTP headers:

\begin{lstlisting}[language=Lisp, caption=Asset Serving Handler]
(defmethod handle-asset-request ((handler asset-handler) request)
  "Handle asset requests with caching headers"
  (let* ((asset-id (get-parameter request "id"))
         (asset-info (get-asset-info *asset-registry* asset-id)))
    (when asset-info
      (let ((response (make-response :ok)))
        (setf (header response "Cache-Control")
              (format nil "max-age=~A" (getf asset-info :max-age)))
        (setf (header response "ETag")
              (getf asset-info :etag))
        (setf (header response "Last-Modified")
              (getf asset-info :last-modified))
        (serve-file response (getf asset-info :path))))))
\end{lstlisting}

\section{Performance Optimization}

\subsection{Asset Preloading}

\subsubsection{Predictive Loading}

The system SHALL implement predictive loading based on:

\begin{itemize}
\item Player movement patterns
\item Scene transition likelihood
\item Asset usage statistics
\item Network conditions
\end{itemize}

\subsubsection{Background Loading}

Non-critical assets SHALL be loaded in background threads:

\begin{itemize}
\item Progressive quality improvement
\item Background texture streaming
\item Ambient sound preloading
\end{itemize}

\subsection{Bandwidth Optimization}

\subsubsection{Adaptive Quality}

Asset quality SHALL adapt based on:

\begin{itemize}
\item Network bandwidth measurements
\item Device capabilities
\item User preferences
\item Battery level (mobile devices)
\end{itemize}

\subsubsection{Compression and Optimization}

Assets SHALL be optimized using:

\begin{itemize}
\item Geometry simplification for distant objects
\item Texture compression (Basis Universal, WebP)
\item Audio compression (Opus)
\item Runtime mesh optimization
\end{itemize}

\section{Security Considerations}

\subsection{Asset Validation}

All assets SHALL be validated to prevent:

\begin{itemize}
\item Malicious asset injection
\item Resource exhaustion attacks
\item Cross-origin resource theft
\item Unauthorized asset access
\end{itemize}

\subsection{Access Control}

Asset access SHALL be controlled through:

\begin{itemize}
\item Authentication requirements
\item Authorization based on user permissions
\item Geographic restrictions
\item Rate limiting for asset requests
\end{itemize}

\section{Error Handling and Fallbacks}

\subsection{Missing Assets}

When assets are unavailable:

\begin{enumerate}
\item Attempt to load lower quality version
\item Use cached version if available
\item Generate procedural fallback
\item Display placeholder geometry
\end{enumerate}

\subsection{Network Failures}

During network interruptions:

\begin{itemize}
\item Serve cached assets when possible
\item Degrade quality gracefully
\item Provide offline mode with limited assets
\item Queue requests for retry when reconnected
\end{itemize}

\section{Testing and Validation}

\subsection{Asset Loading Tests}

\begin{itemize}
\item Load time measurements for different asset types
\item Memory usage analysis
\item Cache hit/miss ratios
\item Network bandwidth utilization
\end{itemize}

\subsection{Integration Tests}

\begin{itemize}
\item Scene loading with multiple assets
\item Progressive loading verification
\item Fallback mechanism testing
\item Cross-platform compatibility
\end{itemize}

\subsection{Performance Benchmarks}

\begin{itemize}
\item Asset loading throughput
\item Memory consumption patterns
\item GPU memory usage
\item Frame rate impact during loading
\end{itemize}

\section{Deployment Strategy}

\subsection{Gradual Rollout}

The extension SHALL be deployed gradually:

\begin{enumerate}
\item \textbf{Phase 1}: Static asset serving with new headers
\item \textbf{Phase 2}: Basic dynamic loading for new scenes
\item \textbf{Phase 3}: Full progressive loading implementation
\item \textbf{Phase 4}: Advanced caching and optimization features
\end{enumerate}

\subsection{Migration Path}

Existing hard-coded assets SHALL be migrated:

\begin{itemize}
\item Convert static references to dynamic loading
\item Generate asset metadata for existing models
\item Update scene descriptions to use new format
\item Maintain backward compatibility during transition
\end{itemize}

\section{Future Extensions}

Potential future enhancements include:

\begin{itemize}
\item Real-time asset updates and hot-reloading
\item Peer-to-peer asset sharing
\item Machine learning-based asset preloading
\item Integration with content delivery networks
\item Advanced compression algorithms
\end{itemize}

\section{References}

\begin{enumerate}
\item Infinity Mode Protocol Specification (Base Protocol)
\item GLTF 2.0 Specification
\item WebGL Asset Loading Best Practices
\item HTTP Caching (RFC 7234)
\item Basis Universal Texture Compression
\end{enumerate}

\end{document}